\begin{scanpage}{0215}{197}
\noindent functions $\Phi$ which enter into Equation (30.04) satisfy relationships (30.05), and may be found by factoring. However, our investigation of the successive waves arising with diffraction at the conductor's ends was approximate. Therefore, it makes no sense to use the precise Expression (30.04) for the primary field. We shall use the approximation expressions for the function $\Phi$
\[\left.\begin{aligned}
\Phi(-k\cos\vartheta,-k\cos\vartheta_0)&=\ln\frac{i}{\gamma ka\sin\dfrac\vartheta2\sin\dfrac{\vartheta_0}2},\\
\Phi(k\cos\vartheta,k\cos\vartheta_0)&=\ln\frac{i}{\gamma ka\cos\dfrac\vartheta2\cos\dfrac{\vartheta_0}2},
\end{aligned}\right\}\tag{32.01}\]
which were obtained by the variational method and have a precision which is sufficient for our purpose (see [83]). More precisely speaking, we will use approximation Equations (32.01) in conjunction with the rigorous Expression (30.05), and we will set
\[\left.\begin{aligned}
\frac1{\Phi(-k\cos\vartheta,-k\cos\vartheta_0)}&=\frac{\displaystyle\ln\frac{i}{\gamma ka\cos\dfrac\vartheta2\cos\dfrac{\vartheta_0}2}}{\displaystyle\ln\frac{2i}{\gamma ka\sin\vartheta}\ln\frac{2i}{\gamma ka\sin\vartheta_0}},\\
\frac1{\Phi(k\cos\vartheta,k\cos\vartheta_0)}&=\frac{\displaystyle\ln\frac{i}{\gamma ka\sin\dfrac\vartheta2\sin\dfrac{\vartheta_0}2}}{\displaystyle\ln\frac{2i}{\gamma ka\sin\vartheta}\ln\frac{2i}{\gamma ka\sin\vartheta_0}}.
\end{aligned}\right\}\tag{32.02}\]
Then the primary field will equal
\[\begin{aligned}
E_\vartheta^{(1)}=H_\varphi^{(1)}={}&-\frac{iE}{2(\cos\vartheta+\cos\vartheta_0)\ln\dfrac{2i}{\gamma ka\sin\vartheta}\ln\dfrac{2i}{\gamma ka\sin\vartheta_0}}\frac{e^{ikR}}{kR}\times\\
&\times\left[\ctg\frac{\vartheta_0}2\ctg\frac\vartheta2\ln\frac{i}{\gamma ka\cos\dfrac\vartheta2\cos\dfrac{\vartheta_0}2}e^{-ikz_1(\cos\vartheta+\cos\vartheta_0)}\right.\\
&\qquad\left.-\tg\frac{\vartheta_0}2\tg\frac\vartheta2\ln\frac{i}{\gamma ka\sin\dfrac\vartheta2\sin\dfrac{\vartheta_0}2}e^{-ikz_2(\cos\vartheta+\cos\vartheta_0)}\right].
\end{aligned}\tag{32.03}\]
Now summing Expressions (31.17) and (32.03), we find the total field scattered by a passive vibrator in the form
\end{scanpage}
